\section{M6 Constructor Anchor Split}
\label{sec:m6_constructor_anchor_split}
\begingroup
\sloppy
\emergencystretch=3em

This fragment refines the global-assembly anchors into constructor-level
handoff tasks from the 20-agent M6 pass.  Each item names the Lean module and
anchors to use, says which explicit hypotheses it removes from downstream
packages, and lists the fields that still have to be supplied by geometry,
partition, or integral-reconstruction workers.

\begin{definition}[Finite active compact boxes]
  \label{anc:m6_compact_active_boxes}
  \lean{Stokes.FiniteActiveOnCompact, Stokes.FiniteActiveOnCompact.ofCompact,
    Stokes.CompactActiveBoxData, Stokes.CompactActiveExtendedBoxData,
    Stokes.CompactActiveExtendedBoxData.toSelectedBoxPartitionOfUnity,
    Stokes.exists_activeCompactCoordinateBoxSelections_piReal}
  \leanok
  \uses{anc:selected_box_partition_of_unity,
    def:partition_subordinate_selected_boxes}
  \emph{Module:} \texttt{Stokes.Global.FiniteActive} and
  \texttt{Stokes.Global.CompactActiveBoxes}.

  \emph{Explicit hypotheses discharged:}
  once a compact control set \(K\) and smooth partition \(\rho\) are fixed,
  \texttt{FiniteActiveOnCompact.ofCompact} supplies the finite active set and
  the pointwise \texttt{fintsupport} subset facts.  Once
  \texttt{CompactActiveBoxData} is built, the per-index lower and upper
  corners, \(a_i \le b_i\), support containment in \(\operatorname{Icc} a_i b_i\),
  and the selected interior box facts are recovered by accessors.  Once
  \texttt{CompactActiveExtendedBoxData} is built, the older
  \texttt{SelectedBoxPartitionOfUnity} fields are obtained without restating
  them index by index.

  \emph{Fields still required:}
  the compact coordinate support set for each active index; compactness of
  those coordinate supports; the equality between the selected
  \texttt{CompactCoordinateBoxSelection.K} and that support; containment of the
  chart representative's \texttt{tsupport} in the coordinate support;
  target-chart and overlap inclusions for every active box; and one open
  smoothness neighborhood with \texttt{ContDiffOn} for every active index.

  \emph{Handoff task:}
  geometry workers should construct \texttt{CompactActiveExtendedBoxData};
  partition workers should consume only
  \texttt{toSelectedBoxPartitionOfUnity} and avoid reopening the box-selection
  fields unless a target/overlap inclusion fails.
\end{definition}

\begin{definition}[Localized support and smoothness controls]
  \label{anc:m6_localized_support_smoothness}
  \lean{Stokes.LocalizedSupportControl,
    Stokes.LocalizedSupportControl.toLocalizedFormData,
    Stokes.LocalizedSupportControl.interiorChartSelectedBox,
    Stokes.CoefficientBoxSupportData,
    Stokes.CoefficientBoxSupportData.toLocalizedSupportControl,
    Stokes.LocalizedSmoothnessData,
    Stokes.LocalizedSmoothnessData.interiorChartExtendedBox}
  \leanok
  \uses{anc:localized_form_data, def:chart_supported_form,
    def:interior_chart_extended_box}
  \emph{Module:} \texttt{Stokes.Global.LocalizedSupport},
  \texttt{Stokes.Global.CoefficientBoxSupport}, and
  \texttt{Stokes.Global.LocalizedSmoothness}.

  \emph{Explicit hypotheses discharged:}
  the raw support calculation for
  \texttt{transitionPullbackInChart} of \(\rho_i \omega\) is replaced by
  \texttt{LocalizedSupportControl}; coefficient-box constructors turn compact
  coordinate-image support data into the exact selected-box support statement;
  \texttt{LocalizedSmoothnessData} packages the open-neighborhood smoothness
  witness needed to upgrade from selected to extended boxes.

  \emph{Fields still required:}
  the concrete source and target charts; lower and upper box corners;
  target and overlap inclusions; a coefficient support or topological-support
  bound for \(\rho_i\); smoothness of the localized chart representative, either
  directly or via chartwise smoothness plus scalar smoothness; and a proof that
  the coefficient being used is the selected partition coefficient.

  \emph{Handoff task:}
  use this layer before building localized interior pieces.  It is the right
  place to solve support and \texttt{ContDiffOn} obligations; downstream
  global constructors should receive a finished support-control or
  extended-box package, not the raw analytic lemmas.
\end{definition}

\begin{definition}[Localized interior pieces]
  \label{anc:m6_localized_interior_pieces}
  \lean{Stokes.LocalizedInteriorPiece,
    Stokes.LocalizedInteriorPiece.localizedFormData,
    Stokes.LocalizedInteriorPiece.extendedBox,
    Stokes.LocalizedInteriorPiece.localStokesData,
    Stokes.LocalizedInteriorPiece.projectLocalEquality,
    Stokes.LocalizedInteriorPieces,
    Stokes.LocalizedInteriorPieces.localProjectEquality,
    Stokes.LocalizedInteriorPieces.sum_projectInterior_eq_artificialBoundary}
  \leanok
  \uses{anc:m6_localized_support_smoothness,
    thm:interior_chart_local_stokes, thm:sum_local_stokes}
  \emph{Module:} \texttt{Stokes.Global.LocalizedInteriorPieces}.

  \emph{Explicit hypotheses discharged:}
  for one active index, \texttt{LocalizedInteriorPiece} turns the coefficient
  \(\rho_i\), support control, and smooth-neighborhood witness into the
  localized form, selected box, extended box, \texttt{InteriorLocalStokesData},
  and project-local interior Stokes equality.  For a finite family,
  \texttt{LocalizedInteriorPieces.localProjectEquality} removes the explicit
  ``prove local Stokes for each active interior piece'' hypothesis, and
  \texttt{sum\_projectInterior\_eq\_artificialBoundary} removes the finite-sum
  local-Stokes algebra.

  \emph{Fields still required:}
  the active finset; the coefficient family; for each index, source chart,
  target chart, lower and upper corners, a \texttt{LocalizedSupportControl},
  and a smooth-neighborhood witness.  This package does not prove that the
  active finset came from a compact partition, that the finite localized sum
  reconstructs \(\omega\), or that artificial faces cancel.

  \emph{Handoff task:}
  interior workers should output \texttt{LocalizedInteriorPieces} plus a map
  from its artificial boundary term to the chosen face-sum model.  The global
  worker can then combine it with \texttt{ArtificialFacePairingData}.
\end{definition}

\begin{definition}[Partition sum-one reconstruction]
  \label{anc:m6_partition_sum_one}
  \lean{Stokes.SmoothPartitionOfUnity.sum_finsupport_univ,
    Stokes.SmoothPartitionOfUnity.sum_eq_one_of_finsupport_subset,
    Stokes.SmoothPartitionOfUnity.sum_eq_one_of_fintsupport_subset,
    Stokes.FiniteActiveOnCompact.coeff_sum_eq_one_on,
    Stokes.FiniteActiveOnCompact.localizedFormSum_apply_eq_self_on,
    Stokes.FiniteActiveOnCompact.localizedFormSum_eqOn,
    Stokes.SelectedBoxPartitionOfUnity.coeff_sum_eq_one_on,
    Stokes.SelectedBoxPartitionOfUnity.localizedFormSum_apply_eq_self_on,
    Stokes.SelectedBoxPartitionOfUnity.localizedFormSum_eqOn}
  \leanok
  \uses{anc:m6_compact_active_boxes, lem:partition_sum_form}
  \emph{Module:} \texttt{Stokes.Global.PartitionSumOne}.

  \emph{Explicit hypotheses discharged:}
  the ad hoc hypothesis
  \[
    \forall x \in K,\quad \sum_{i \in s} \rho_i(x)=1
  \]
  is now obtained from \texttt{FiniteActiveOnCompact} or
  \texttt{SelectedBoxPartitionOfUnity}.  The pointwise reconstruction of the
  localized finite form sum on \(K\) is also available directly, so
  reconstruction workers do not need to manipulate mathlib's
  \texttt{finsupport} and \texttt{fintsupport} APIs.

  \emph{Fields still required:}
  the compact set \(K\) must actually contain the support region on which the
  global integral is reconstructed; the coefficient family in later
  constructors must be definitionally or propositionally the stored partition;
  and any theorem needing equality outside \(K\) must either enlarge \(K\) or
  supply an everywhere coefficient-sum hypothesis separately.

  \emph{Handoff task:}
  use this anchor to fill the \texttt{lem:partition\_sum\_form} side of the
  blueprint.  For exterior-derivative reconstruction, pass the resulting
  coefficient-sum fact to the \texttt{ExtDerivPartitionReconstructionData}
  constructors only after checking the statement is on the same domain.
\end{definition}

\begin{definition}[Artificial face pairing]
  \label{anc:m6_artificial_face_pairing}
  \lean{Stokes.ArtificialFaceIndex, Stokes.artificialFaceIndexSet,
    Stokes.artificialFaceTerm, Stokes.interiorBoundaryFace_sum_eq_indexSet_sum,
    Stokes.interiorBoundaryFace_sum_eq_zero_of_pairing,
    Stokes.ArtificialFacePairingData,
    Stokes.ArtificialFacePairingData.toArtificialBoundaryCancellationData,
    Stokes.ArtificialFacePairingData.toArtificialBoundaryCancellationData_of_boundaryTerm,
    Stokes.GlobalStokesData.interiorBoundaryCancellation_of_facePairing}
  \leanok
  \uses{thm:artificial_faces_vanish, thm:global_compact_support_stokes}
  \emph{Module:} \texttt{Stokes.Global.ArtificialFacePairing}.

  \emph{Explicit hypotheses discharged:}
  the single final-field proof
  \[
    \sum_x \sum_p \operatorname{interiorBoundaryTerm}(x,p)=0
  \]
  can now be supplied by a face-level involution.  The module flattens the
  chart/piece/face indices, proves that the nested sum equals the flattened
  sum, and converts the pairing package into
  \texttt{ArtificialBoundaryCancellationData} or directly into the
  \texttt{GlobalStokesData.interiorBoundaryCancellation} shape.

  \emph{Fields still required:}
  active charts; interior pieces per chart; face labels per interior piece; the
  signed face-term function; a set-preserving pairing on flattened face
  indices; the involution proof; opposite-term cancellation for paired faces;
  zero fixed-point terms; and, when an external per-piece artificial boundary
  term is already chosen, an equality identifying it with the sum of recorded
  face terms on active pieces.

  \emph{Handoff task:}
  local-geometry workers should provide the pairing and sign equations.
  Assembly workers should consume only
  \texttt{toArtificialBoundaryCancellationData} or
  \texttt{interiorBoundaryCancellation\_of\_facePairing}; they should not expand
  the flattened sigma index set unless debugging signs.
\end{definition}

\begin{definition}[Boundary chart-change pieces]
  \label{anc:m6_boundary_chart_change_pieces}
  \lean{Stokes.BoundaryChartChangePieceData,
    Stokes.BoundaryChartChangeSelectedPieceData,
    Stokes.BoundaryChartChangeExtendedPieceData,
    Stokes.BoundaryChartChangeFamilyData,
    Stokes.BoundaryChartChangeSelectedFamilyData,
    Stokes.BoundaryChartChangeExtendedFamilyData,
    Stokes.BoundaryChartChangeFamilyData.toChartChangeCancellationData_selected,
    Stokes.BoundaryChartChangeFamilyData.toChartChangeCancellationData_extended,
    Stokes.ProjectLocalGlobalStokesData.chartChangeCancellation_of_boundaryChartChange_selected,
    Stokes.ProjectLocalGlobalStokesData.chartChangeCancellation_of_boundaryChartChange_extended}
  \leanok
  \uses{thm:boundary_chart_integral_invariance,
    anc:project_local_global_stokes}
  \emph{Module:} \texttt{Stokes.Global.BoundaryChartChangePieces}.

  \emph{Explicit hypotheses discharged:}
  the finite-sum field
  \texttt{ProjectLocalGlobalStokesData.chartChangeCancellation} no longer has
  to be proved by hand.  A per-piece oriented change-of-variables package plus
  a target selected or extended boundary box gives pointwise equality from the
  source project-local boundary term to the chosen boundary-partition term; the
  family wrapper then sums those equalities in the exact final-field shape.

  \emph{Fields still required:}
  for every active boundary piece: target boundary chart, target lower and
  upper corners, \texttt{boundaryChartOrientedChangeOfVariables}, and either a
  selected or extended target boundary box.  At the family level, one still
  needs the equality identifying \texttt{D.boundaryPartitionTerm} with the
  transported target project-local boundary integral, plus the
  \texttt{IsManifold I 1 M} instance required by the chart-change wrappers.

  \emph{Handoff task:}
  boundary workers should build a selected-family or extended-family package
  immediately after constructing image data.  Global workers can then fill
  \texttt{chartChangeCancellation} by one theorem call.
\end{definition}

\begin{definition}[Boundary-piece global constructors]
  \label{anc:m6_boundary_piece_constructors}
  \lean{Stokes.BoundaryProjectLocalPieces,
    Stokes.BoundaryProjectLocalPieces.toProjectLocalGlobalStokesData,
    Stokes.BoundaryProjectLocalPieces.toGlobalStokesData,
    Stokes.BoundaryProjectLocalPieces.stokes,
    Stokes.OrientedBoundaryProjectLocalPieces,
    Stokes.OrientedBoundaryProjectLocalPieces.toGlobalStokesData_of_orientedAtlas,
    Stokes.OrientedBoundaryProjectLocalPieces.toGlobalStokesData_of_orientedManifold,
    Stokes.OrientedBoundaryProjectLocalPieces.stokes_of_orientedAtlas,
    Stokes.OrientedBoundaryProjectLocalPieces.stokes_of_orientedManifold}
  \leanok
  \uses{anc:m6_boundary_chart_change_pieces,
    thm:boundary_chart_local_stokes, anc:project_local_global_stokes}
  \emph{Module:} \texttt{Stokes.Global.BoundaryPieces}.

  \emph{Explicit hypotheses discharged:}
  \texttt{BoundaryProjectLocalPieces} derives local project Stokes from stored
  extended boundary-chart boxes and instantiates
  \texttt{ProjectLocalGlobalStokesData}.  The oriented variant uses source
  extended boxes, target selected boxes, and image data to instantiate
  \texttt{GlobalStokesData} directly with no interior pieces.

  \emph{Fields still required:}
  active charts and local pieces; source and target chart choices; source box
  corners; extended source boundary boxes; the global bulk and boundary real
  numbers; bulk reconstruction; boundary-partition reconstruction; and the
  chart-change cancellation field.  The oriented variant additionally needs
  target chart choices, target corners, target selected boxes, image data, and
  either oriented-atlas chart membership hypotheses or a
  \texttt{BoundaryChartOrientedManifold} instance.

  \emph{Handoff task:}
  use this constructor for boundary-only tests and for validating the boundary
  side of mixed global Stokes before interior artificial-face cancellation is
  connected.
\end{definition}

\begin{definition}[Interior and mixed global constructors]
  \label{anc:m6_interior_mixed_constructors}
  \lean{Stokes.InteriorProjectLocalPieces,
    Stokes.InteriorProjectLocalPieces.ofCancellationData,
    Stokes.InteriorProjectLocalPieces.toGlobalStokesData,
    Stokes.interiorProjectLocalGlobalStokes,
    Stokes.MixedInteriorPackage, Stokes.MixedBoundaryPackage,
    Stokes.MixedGlobalStokesData,
    Stokes.MixedGlobalStokesData.toGlobalStokesData,
    Stokes.mixedGlobalStokes}
  \leanok
  \uses{anc:m6_localized_interior_pieces,
    anc:m6_artificial_face_pairing, anc:m6_boundary_piece_constructors,
    anc:global_stokes_data}
  \emph{Module:} \texttt{Stokes.Global.InteriorGlobalConstructor} and
  \texttt{Stokes.Global.MixedGlobalConstructor}.

  \emph{Explicit hypotheses discharged:}
  the interior-only constructor turns stored \texttt{InteriorLocalStokesData}
  packages plus artificial-boundary cancellation into a final
  \texttt{GlobalStokesData} instance with empty genuine boundary pieces.  The
  mixed constructor separates the final theorem into three handoff outputs:
  reconstruction data, an interior local-Stokes package, and a boundary
  local-Stokes package.  Its \texttt{toGlobalStokesData} definition performs
  only record alignment.

  \emph{Fields still required:}
  for the interior-only path: active charts, interior pieces, local Stokes data
  per piece, global bulk reconstruction, artificial-boundary cancellation, and
  the proof that the represented boundary integral is zero.  For the mixed
  path: a \texttt{PartitionReconstructionData} package, interior and boundary
  boundary-term functions, local Stokes packages for both sides, artificial
  interior-boundary cancellation, and boundary chart-change cancellation.

  \emph{Handoff task:}
  use the interior-only constructor as a narrow smoke test for
  \texttt{LocalizedInteriorPieces} plus \texttt{ArtificialFacePairingData}.
  Use \texttt{MixedGlobalStokesData} as the main parent-agent assembly target.
\end{definition}

\begin{definition}[Partition and exterior-derivative reconstruction constructors]
  \label{anc:m6_reconstruction_constructors}
  \lean{Stokes.localizedFormSum,
    Stokes.localizedFormSum_apply_eq_self_on,
    Stokes.PartitionReconstructionData,
    Stokes.PartitionReconstructionData.toGlobalStokesData,
    Stokes.extDeriv_finset_sum,
    Stokes.extDerivWithin_finset_sum,
    Stokes.extDeriv_transitionPullbackInChart_localizedFormSum,
    Stokes.ExtDerivPartitionReconstructionData,
    Stokes.ExtDerivPartitionReconstructionData.ofPartitionReconstructionData,
    Stokes.ExtDerivPartitionReconstructionData.ofPartitionReconstructionData_of_coeff_sum_eq_one,
    Stokes.ExtDerivPartitionReconstructionData.toGlobalStokesData}
  \leanok
  \uses{anc:m6_partition_sum_one, lem:extderiv_finite_sum,
    lem:partition_sum_form, thm:global_bulk_eq_local_sum,
    thm:global_boundary_eq_local_sum}
  \emph{Module:} \texttt{Stokes.Global.Reconstruction} and
  \texttt{Stokes.Global.ExtDerivReconstruction}.

  \emph{Explicit hypotheses discharged:}
  \texttt{PartitionReconstructionData} isolates the two global reconstruction
  equalities and converts them, together with local/cancellation/chart-change
  inputs, into the final \texttt{GlobalStokesData} record.  The exterior
  derivative module supplies the finite-sum \texttt{extDeriv} wrappers and an
  augmented reconstruction package whose chartwise exterior-derivative
  equality can be carried as a named field or produced from an everywhere
  coefficient-sum-one hypothesis.

  \emph{Fields still required:}
  finite active chart labels; interior and boundary piece finsets; local bulk
  term functions; boundary partition terms; represented global bulk and
  boundary values; the two reconstruction equalities; local Stokes on all
  pieces; artificial-face cancellation; chart-change cancellation; and, for
  the exterior-derivative package, the chartwise equality comparing
  \texttt{extDeriv} of the localized finite sum with \texttt{extDeriv} of the
  original chart representative.

  \emph{Handoff task:}
  reconstruction workers should build this package before calling any final
  Stokes theorem.  Analytic workers should decide whether the ext-derivative
  field is proved from \texttt{PartitionSumOne} or recorded as an explicit
  chartwise theorem.
\end{definition}

\begin{theorem}[M6 constructor dependency order]
  \label{anc:m6_constructor_dependency_order}
  \lean{Stokes.mixedGlobalStokes, Stokes.globalStokes,
    Stokes.projectLocalGlobalStokes}
  \leanok
  \uses{anc:m6_compact_active_boxes,
    anc:m6_localized_support_smoothness,
    anc:m6_localized_interior_pieces,
    anc:m6_partition_sum_one,
    anc:m6_artificial_face_pairing,
    anc:m6_boundary_chart_change_pieces,
    anc:m6_reconstruction_constructors}
  \emph{Explicit hypotheses discharged:}
  this order discharges no mathematical field by itself; it records which
  constructor should be used before a parent worker asks for the next final-data
  equality.

  \emph{Fields still required:}
  the union of the fields listed in the preceding constructor anchors.  In
  practice, missing fields should be assigned to the earliest constructor that
  names them, not patched into \texttt{GlobalStokesData} directly.

  The intended parent-agent integration order is:
  first build compact active boxes and partition sum-one reconstruction; then
  construct localized interior pieces from support and smoothness controls;
  then provide artificial face pairings for their boundary terms; then build
  boundary chart-change family data for genuine boundary pieces; finally feed
  the reconstruction, local-Stokes, cancellation, and chart-change packages
  into \texttt{MixedGlobalStokesData} or directly into
  \texttt{GlobalStokesData}.
\end{theorem}

\endgroup
